\section*{Chapter \ref{ch:g10:economic-models} --- Economic Models}
% figdata: figdata/grade-10/01-economic-models/out/ (trend_points, trend_fit, trend_misses,
% yield_groups, brenland_gdppc); every computation follows the rounded values printed.

\begin{solution}{exo:g10:economic-models:1}
In 1929, about \usd{12000} in the United States and about \usd{7000} in Argentina; in 2022, about
\usd{58000} and about \usd{18000} (the table behind the chart gives \usd{11954}, \usd{6961},
\usd{58488} and \usd{18292}). The American curve stayed much closer to its line: it ends a little
above it, while the Argentine curve ends far below.
\end{solution}

\begin{solution}{exo:g10:economic-models:2}
1.~Positive (and checkable in the Maddison data). 2.~Normative (``ought''). 3.~Positive: a
prediction, which data on free-fare experiments could test. 4.~Normative (``unfair''). 5.~Positive.
\end{solution}

\begin{solution}{exo:g10:economic-models:3}
In $y_t = y_0(1+g)^t$, $y_t$ is endogenous; $y_0$, $g$ and $t$ are exogenous. In $R = p\,q$, the
takings $R$ are endogenous; $p$ and $q$ are exogenous.
\end{solution}

\begin{solution}{exo:g10:economic-models:4}
Four: Canada (87\,\%), Chile (88\,\%), the United States (117\,\%) and Switzerland (122\,\%).
Argentina ended furthest below (56\,\%), Brazil furthest above (987\,\%, almost ten times its line's
value).
\end{solution}

\begin{solution}{exo:g10:economic-models:5}
No. The statement holds \omterm{def:g10:economic-models:other-things-equal}{other things equal}: it compares two situations that differ only by the price.
Between May and July other things changed --- the weather is warmer, there are more visitors --- and
these raised the number of ice creams people wanted at every price. The observation mixes the effect
of the price with the effect of summer, so it cannot contradict a statement about the price alone.
\end{solution}

\begin{solution}{exo:g10:economic-models:6}
After 10 years: $20000 \times 1.02^{10} \approx 20000 \times 1.2190 \approx \usd{24380}$. After 35
years: $20000 \times 1.02^{35} \approx 20000 \times 1.9999 \approx \usd{39998}$. So income doubles in
about 35 years (the rule of 70 gives $70/2 = 35$).
\end{solution}

\begin{solution}{exo:g10:economic-models:7}
The ratio is $8074/2702 \approx 2.988$; the yearly factor is $2.988^{1/59} \approx 1.0187$, a rate of
1.87\,\% a year. Forecast for 2022: $8074 \times 1.0187^{93} \approx 8074 \times 5.6015 \approx
\usd{45227}$. The observed value, \usd{45530}, is $45530/45227 \approx 1.007$ times the forecast:
the forecast was only 0.7\,\% too low. The chart of misses puts Canada at 87\,\% instead: the line through
all sixty points of 1870--1929 grows at 2.01\,\% a year, faster than the two-point 1.87\,\% (Canada's
1870 value lies above that line and its 1929 value below it), so it predicts more for 2022 and the observation falls short of it.
The method of fitting matters, as the measure box of the chapter says.
% figdata: trend_misses (CAN rate 2.01, pct 87)
\end{solution}

\begin{solution}{exo:g10:economic-models:8}
The ratio is multiplied each year by $1.02/1.01 \approx 1.0099$, so after 70 years by $1.0099^{70}
\approx 1.99$: the faster country is about twice as rich. (Directly: $1.02^{70} \approx 4.00$ and
$1.01^{70} \approx 2.01$.) On the opening chart the two rates differ much less. The ratio of
Argentine to American income is multiplied each year by $1.0164/1.0153 \approx 1.0011$, so over 93
years by $1.0011^{93} \approx 1.11$. The model predicts a ratio of $0.582 \times 1.11 \approx 0.65$
in 2022: Argentina at about 65\,\% of the American level, a little higher than in 1929.
\end{solution}

\begin{solution}{exo:g10:economic-models:9}
$4900/920 \approx 5.3$: more than five times. It is not the \omterm{def:g10:economic-models:causal-effect}{causal effect} because the countries that
use much fertiliser also differ in other ways that raise yields (machines, irrigation, seed, roads:
\omterm{def:g10:economic-models:confounder}{confounding factors}, often going with a higher income), and because good land and rain make
fertiliser more worth buying (\omterm{def:g10:economic-models:reverse}{reverse causation}). The ratio mixes all of these.
\end{solution}

\begin{solution}{exo:g10:economic-models:10}
Income per person (or the wealth of the country). Richer countries have more phones per person, since
people can afford them; they also have more doctors, cleaner water and better food, which raise life
expectancy. So phones and life expectancy rise together across countries even if phones themselves
added nothing to life expectancy.
\end{solution}

\begin{solution}{exo:g10:economic-models:11}
Firms with large sales have more money to spend and more to gain from advertising (each extra buyer
reached is worth more when the firm already sells widely), and many firms set their advertising
budget as a share of their sales. So large sales cause large advertising budgets. An observation that
points to an effect of advertising: sales rising in the towns where a campaign was run, chosen by
lot, and not in comparable towns where it was not; or sales jumping right after a campaign started
for reasons unrelated to the season.
\end{solution}

\begin{solution}{exo:g10:economic-models:12}
For example: (1) the growth rate of one decade very different from that of the next (2\,\% in one
decade, 5\,\% in the next), or a fall over several years: this contradicts the assumption that the
rate is the same every year. (2) After a deep fall (a war, a depression), income never coming back to
the old line but growing parallel to it, below it: the model assumes that nothing interrupts the
path, so a lasting step down contradicts it. (A curve whose slope on a times-ten scale keeps falling,
decade after decade, also contradicts the constant rate.)
\end{solution}

\begin{solution}{exo:g10:economic-models:13}
Draw lots among the 200 farms: for example, 100 names drawn from a box receive \qty{50}{kg} of
fertiliser per hectare, and the other 100 receive none (or keep their usual practice). At harvest,
measure the maize yield of each farm in kilograms per hectare. The difference between the average
yield of the treated farms and that of the others estimates the effect of the fertiliser
(\cref{prop:g10:economic-models:lottery}). Volunteers would differ from the others: farmers with better
land, more skill or more money are more likely to ask for fertiliser, and they would have had higher
yields anyway. The comparison of volunteers with the others would then most likely overstate the
effect.
\end{solution}

\begin{solution}{exo:g10:economic-models:14}
Ratio in 1929: $11954/6961 \approx 1.717$. Predicted for 2022: $50167/32520 \approx 1.543$. Observed in
2022: $58488/18292 \approx 3.197$. The observed ratio is $3.197/1.543 \approx 2.07$ times the
predicted one. The United States ended $58488/50167 \approx 1.166$ times its forecast, Argentina
$18292/32520 \approx 0.562$ times its forecast; the ratio's error is $1.166/0.562 \approx 2.07$, and
the Argentine part (a factor $1/0.562 \approx 1.78$) is much larger than the American part ($1.166$).
The model predicted the ratio would fall (\cref{prop:g10:economic-models:ratio}, since Argentina's line
grew faster); instead it almost doubled. Both countries grew, but the model's prediction is about the
rates staying constant, and a ratio that rose instead of falling shows that the two rates did not stay
as they were.
\end{solution}

\begin{solution}{exo:g10:economic-models:15}
The researcher must assume that, without the flood, the incomes of eastern and western families would
have been the same in 2020 on average --- or, more weakly, would have changed in the same way. Data from
before 2019 help: if eastern and western incomes were similar and moved together in 2010--2018, the
assumption is more believable, and the comparison can be made on the change since 2018 rather than the
level. If eastern incomes were already lower, or already falling relative to the west (poorer soil,
fewer roads), the comparison would mix that difference with the flood, and the assumption is in doubt.
\end{solution}

\begin{solution}{pb:g10:economic-models:1}
\textbf{1.} $27027/24675 \approx 1.0953$: a rise of 9.53\,\%.

\textbf{2.} $25840/27027 \approx 0.9561$: a fall of 4.39\,\%.

\textbf{3.} $28371/29997 \approx 0.9458$: a fall of 5.42\,\%.

\textbf{4.} Real GDP per person of 2025 is endogenous: the model computes it. The growth rate is
exogenous: the model takes it as given (here, from the past).

\textbf{5.} Normative: ``too slowly'' compares with what ought to have happened. A positive version:
``Real GDP per person grew by less from 2007 to 2014 than from 2000 to 2007'' ($27905/27027 \approx
1.032$ against $1.095$), which the table confirms.

\textbf{6.} The ratio is $1.0953$ (question 1). The yearly factor is $1.0953^{1/7} \approx 1.0131$, so
$g = 1.31\,\%$ a year.

\textbf{7.} $24675 \times 1.0131^{7} \approx 24675 \times 1.0954 \approx 27029$, against \dal{27027}
observed: the fit is exact up to rounding, as it must be with two points.

\textbf{8.} About $70/1.31 \approx 53$ years (check: $1.0131^{53} \approx 1.99$).

\textbf{9.} That real GDP per person keeps growing at 1.31\,\% a year, with no interruption. A
recession, a war, a large natural disaster, or a lasting change in the growth rate (a new technology
spreading, an ageing workforce) would break it.

\textbf{10.} $27027 \times 1.0131^{2} \approx 27027 \times 1.0264 \approx \text{\dal{27741}}$. The observed
value is $25840/27741 \approx 0.931$ times the prediction: 6.9\,\% below.

\textbf{11.} $27027 \times 1.0131^{18} \approx 27027 \times 1.2640 \approx \text{\dal{34162}}$.

\textbf{12.} $31470 - 34162 = -2692$, and $-2692/34162 \approx -0.079$: a forecast error of
$-7.9\,\%$. The forecast was too high by 7.9\,\% of its value.

\textbf{13.} The ratio of forecast to observation is $34162/31470 \approx 1.086$, and $1.0131^{6}
\approx 1.081$: the gap is about six years of growth at the rate $g$. Brenland reached in 2025 the
level the model had expected for about 2019.

\textbf{14.} $29997/26170 \approx 1.1462$, and $1.1462^{1/9} \approx 1.0153$: 1.53\,\% a year, a little
faster than $g = 1.31\,\%$. Growth did not slow after the recession; if anything it was faster for a
decade, as part of the lost ground was made up.

\textbf{15.} $25840 \times 1.0131^{16} \approx 25840 \times 1.2315 \approx \text{\dal{31822}}$. Error:
$(31470 - 31822)/31822 \approx -0.011$, that is $-1.1\,\%$.

\textbf{16.} Restarting from 2009 with the same rate cuts the error from $-7.9\,\%$ to $-1.1\,\%$.
The rate held roughly; what failed is the assumption that nothing interrupts the path. The recession
of 2008--2009 lowered the level, and the later years grew from the lower level instead of returning
to the old line.

\textbf{17.} No: a model of a constant yearly rate is about the path over many years, and one bad
year followed by a quick return is a small departure from it. What would contradict it is a pattern:
several years in a row far from the rate, or a fall from which the series never comes back to the
line, as after 2009.

\textbf{18.} $27905/24675 \approx 1.1309$, and $1.1309^{1/14} \approx 1.0088$: 0.88\,\% a year.
Forecast: $27905 \times 1.0088^{11} \approx 27905 \times 1.1012 \approx \text{\dal{30729}}$. Error:
$(31470 - 30729)/30729 \approx +0.024$, that is $+2.4\,\%$ (too low).

\textbf{19.} No. The 2014 forecaster had seven more years of data and started closer to 2025 (11 years
instead of 18), so smaller errors are expected. And the fitting period 2000--2014 happens to contain
the recession, which pulls the rate down: the forecast was too low, the 2007 forecast too high, and
neither method was tested on more than one case. A method is judged on many forecasts, not on one.

\textbf{20.} The constant-growth forecast made in 2007 overstated Brenland's real GDP per person in
2025 by 7.9\,\% of its value (\dal{34162} forecast against \dal{31470} observed); almost all of the
error is due to the recession of 2008--2009, since the same rate restarted from 2009 misses by only
1.1\,\%.
\end{solution}
